\documentclass[10pt]{article}

% Specify the margins and text size.
\setlength{\textwidth}{6.4in}
\setlength{\textheight}{9.5in}
\setlength{\oddsidemargin}{0pt}
\setlength{\evensidemargin}{0pt}
\setlength{\topmargin}{0pt}
\setlength{\hoffset}{.05in}
\setlength{\voffset}{-1in}

\setlength{\parskip}{5pt}
\setlength{\parindent}{0pt}

% Load some fonts and symbol packages
\usepackage{latexsym}
\usepackage{pifont}       % contains 'star' symbol for counterinsurgency handbook title
\usepackage{yfonts} 
\usepackage{amsmath}
\usepackage{amsfonts}

\usepackage{graphicx}     % actually, this is loaded by pstricks
\usepackage[T1]{fontenc}
\usepackage{ifthen}
\usepackage{pstricks,pst-grad,pst-text,pst-node,multido,pst-plot,calc,pst-3dplot}
%\usepackage[all]{xy}
%\usepackage{animate}

% The hyperref package inserts links into the pdf.
\definecolor{MyLinkColor}{rgb}{.1,.2,1}
\definecolor{MyCiteColor}{rgb}{.1,1,.2}
\definecolor{MyURLColor}{rgb}{.4,.4,.4}
\usepackage[backref=true,pagebackref=false,hyperindex,colorlinks=true,
  linkcolor=MyLinkColor,urlcolor=MyURLColor]{hyperref}


% The tweaklist package is something I found on the web.  It provides a simple interface
% for making changes to spacing used in the itemize and enumerate environments.  Comment
% this out if you don't care to use tweaklist.
\usepackage{tweaklist}
\renewcommand{\itemhook}{\setlength{\parskip}{2pt}\setlength{\parsep}%
{1pt}\setlength{\topsep}{0pt}\setlength{\itemsep}{0pt}}

\newcommand{\U}{\underline{\hspace{5pt}}}

\usepackage{listings}
\newcommand{\Z}{\hphantom{0}}

\begin{document}
\pagestyle{empty}
\lstset{language=R, showspaces=false, showstringspaces=false}

\href{http://www.su.edu}{\includegraphics[height=1.75cm]{sulogo.eps}}
\vspace{-1.69cm}

{{\ }\hfill\small
\begin{tabular}{cl}
& Math 207\\
& Introduction to Statistics\\
\end{tabular}
}
\setlength{\baselineskip}{1.05\baselineskip}
\bigskip

\begin{center}
\textbf{\large  Quiz 2}
\end{center}
\medskip


1. (a) For each list below, calculate the mean and the sd.

\hspace{20pt} (i) 1, 3, 4, 5, 7\vspace{1.8in}

\hspace{20pt} (ii) 6, 8, 9, 10, 12\vspace{1.8in}

\hphantom{1. } (b) How is list (ii) related to list (i)?  How does this
relationship carry over to the mean?  The sd?
\vspace{1.25in}

2. Convert the values in 1(a)(i) to standard units.

\vfill
\eject
{\ }

3. (a) A list of numbers has $\mbox{mean}=4$ and $\mbox{sd}=2$.
Convert $x=6$ to standard units.
\vspace{1.25in}

\hphantom{2. } (b) What  value $x$ has value $-2$ in standard units?  
Assume that the mean is 4 and sd is 2.
\vspace{1.5in}

4. If $x=5$, then $z=1$.  If $z=-2$, then $x=-4$.  Find the mean and sd.
\vspace{2.5in}

5. Compute the RMS for each list.

\hspace{10pt} a) $1$, $-1$, $1$, $-1$
\vspace{.8in}

\hspace{10pt} b) $1$, $1$, $1$, $1$, $1$
\vspace{.8in}

\hspace{10pt} c) $2$, $-3$, $1$, $-1$, $7$, $-4$
\vfill
\eject

\end{document}

